%=============================================================================!
% 7.0  CHAPTER OPENING                                                        !
%=============================================================================!
Chapter~6 described motion through generalised coordinates and their
velocities, and introduced the conjugate momentum belonging to each
coordinate. We now use those momenta in place of the velocities. This
change leads to Hamilton's equations and makes the state of a system a
point in phase space, the space of all allowed coordinates and momenta.
Following many starting states at once reveals properties of the motion
that are difficult to see from a single trajectory.

For the time-independent systems considered here, the flow preserves
energy as well as the area, or in more dimensions the volume, of a region
of states. We derive this volume conservation as Liouville's theorem,
which Chapter~12 needs for statistical mechanics. We then connect spatial
symmetries to conserved momenta and establish the conditions under which
reversing momenta retraces a trajectory. These results give us concrete
tests for a numerical method: the final section compares two simple step
maps, preparing the integration methods of Chapter~9.
